\chapter{\emph{Delta WCOJ Propagation}}
\label{cap:delta-wcoj}

\begin{intuicion}
La técnica que el motor adoptará en stage 2b para resolver el
incremental update de k-hop reachability viene de
Kankanamge-Sahu-Mhedhbi-Salihoglu, ``Graphflow: An Active Graph
Database'' (SIGMOD 2017). La idea de fondo: no recomputar todo el
índice, sino propagar el \emph{delta} (la diferencia entre el
estado nuevo y el viejo) hacia los predecesores hasta que el delta
se vacíe o se alcance \texttt{MAX\_DEPTH}. Es una versión
``targeted'' del BFS reverse multi-source: en lugar de visitar
todos los vértices, sólo visita los que la nueva arista
genuinamente afecta. Este capítulo pseudo-codifica el algoritmo,
prueba soundness, y discute el cascade budget que protege contra
hubs patológicos.
\end{intuicion}

\section{El algoritmo}

\begin{algorithm}[H]
\caption{\texttt{khop\_reach.delta\_update(u, v, t)} — esquema}
\KwIn{nueva arista \((u, v, t)\), grafo \(G\), índice \texttt{khop\_reach}}
\BlankLine
\Comment{Round 1 — efecto directo: $u$ ahora alcanza $\mathrm{type}(v)$ y todo lo que $v$ alcanza.}
\(\Delta_u \leftarrow \bits[v] \setminus \bits[u]\)\;
\If{\(\Delta_u = 0\) \textbf{and} \(t\) no mejora $\mathrm{min/max\_first}[u]$}{
    \Return\,(la arista no cambia el índice)\;
}
\(\bits[u] \mathrel{|}\!= \bits[v]\)\;
\For{cada \(L\) en \(\bits[v]\)}{
    \(\mathrm{min\_first}[u][L] \leftarrow \min(\mathrm{min\_first}[u][L], t)\)\;
    \(\mathrm{max\_first}[u][L] \leftarrow \max(\mathrm{max\_first}[u][L], t)\)\;
}
\BlankLine
\Comment{Rounds 2..MAX\_DEPTH — propagación reverse.}
\texttt{frontier} \(\leftarrow \{u\}\)\;
\texttt{visited} \(\leftarrow \{u\}\)\;
\For{$d = 2$ \textbf{to} \(\texttt{MAX\_DEPTH}\)}{
    \texttt{next} \(\leftarrow \emptyset\)\;
    \For{\(x \in\) \texttt{frontier}}{
        \For{\(w \in N^{-}(x)\) (\emph{reverse adjacency})}{
            \(\Delta_w \leftarrow \bits[x] \setminus \bits[w]\)\;
            \If{\(\Delta_w \neq 0\) \textbf{or} \(t_{wx}\) mejora \(\mathrm{min/max\_first}[w]\)}{
                \(\bits[w] \mathrel{|}\!= \bits[x]\)\;
                actualizar \(\mathrm{min/max\_first}[w][L]\) para cada \(L \in \bits[x]\), usando \(t_{wx}\) (timestamp de la arista \(w \to x\))\;
                \If{\(w \notin\) \texttt{visited}}{
                    \texttt{next} \(\leftarrow\) \texttt{next} \(\cup \{w\}\)\;
                    \texttt{visited} \(\leftarrow\) \texttt{visited} \(\cup \{w\}\)\;
                }
            }
        }
    }
    \texttt{frontier} \(\leftarrow\) \texttt{next}\;
    \If{\(|\texttt{visited}| > \texttt{CASCADE\_BUDGET}\)}{
        \emph{abort}: invalidar \texttt{khop\_reach}; rebuild en próximo hunt\;
        \Return\;
    }
}
\end{algorithm}

\section{Análisis paso a paso}

\subsection{Round 1: efecto directo}

La nueva arista \(u \to v\) significa que \(u\) ahora puede
alcanzar:
\begin{itemize}
    \item \(\mathrm{type}(v)\) directamente (1 salto).
    \item Todo lo que \(v\) ya alcanzaba (vía \(v\)).
\end{itemize}

\noindent
Eso es exactamente \(\bits[u] \mathrel{|}\!= \bits[v]\) (incluyendo
el bit de \(\mathrm{type}(v)\) si estaba en \(\bits[v]\) por el
caso degenerate ``\(v\) alcanza su propio tipo''; en la
implementación real se asegura prendido).

Si \(\bits[v] \setminus \bits[u] = 0\) (\(\Delta_u\) vacío), \(u\)
no gana ningún tipo nuevo. Pero todavía puede ganar mejor
\(\mathrm{min/max\_first}\) (si \(t\) es más temprano que cualquier
camino previo) ---chequeo separado.

\subsection{Rounds 2..MAX\_DEPTH: propagación reverse}

Para cada \(w\) que tiene una arista directa a \(u\), recursivamente
preguntamos si \(w\) gana algo. Si \(\Delta_w = 0\) y los timestamps
no mejoran, se aborta esta rama (no se agrega a la frontier).

La cota \texttt{MAX\_DEPTH} en \(d\): después de
\texttt{MAX\_DEPTH} rounds, la propagación se detiene. Razón: el
índice sólo guarda alcance \(\leq \texttt{MAX\_DEPTH}\) saltos;
un cambio más profundo no afecta a la sobre-aproximación
guardada.

\subsection{Cascade budget}

Si \texttt{visited} pasa un umbral (default 100\,000), abortar la
propagación e invalidar el índice completo. El próximo
\texttt{ensure\_khop\_reach} reconstruye desde cero
(capítulo~\ref{cap:khop-detalle}).

¿Por qué? Porque los grafos scale-free tienen hubs raros donde una
cascada visita millones. En esos casos el rebuild \(O(|\E|)\) es
\emph{cheaper} que la cascada \(O(\)visited\()\). Mejor caer al
rebuild que pagar más en propagación.

\section{Soundness}

\begin{theorem}[Delta update preserva sobre-aproximación]
Sea \(G_t\) el grafo después de \(t\) llamadas a
\texttt{add\_relation}, e índice \(I_t\) computed por aplicación
secuencial de \texttt{delta\_update}. Entonces para todo \(v\) y
\(L\):
\[
I_t.\bits[v][L] = 1 \implies \exists\,\text{walk}\;v
\rightsquigarrow_{\leq \texttt{MAX\_DEPTH}} L\;\text{en}\;G_t.
\]
Es decir: el índice puede sobre-aproximar (decir 1 cuando no
existe walk) pero NUNCA puede sub-aproximar.
\end{theorem}

\begin{proof}[Idea]
Inducción sobre \(t\). Base \(t=0\): \(I_0\) es construido por el
builder (capítulo~\ref{cap:khop-detalle}), que es exacto. Paso
inductivo: \texttt{delta\_update} sólo \emph{agrega} bits via
\(\bits[w] \mathrel{|}\!=\); nunca quita bits. Y sólo
\emph{baja} \(\mathrm{min\_first}\) y \emph{sube}
\(\mathrm{max\_first}\); nunca al revés. Por lo tanto, cualquier
``walk válido en \(G_t\) que el índice marca'' sigue siendo válido
en \(G_{t+1}\) (porque \(G_{t+1} \supseteq G_t\)). Sound. \(\square\)
\end{proof}

\subsection{Falsos positivos OK; falsos negativos NO}

La sobre-aproximación es OK porque el matcher \emph{verifica}
después: si el prune deja pasar un start que no debería, el DFS
lo descarta sin emit. Costo: un poco más de DFS.

La sub-aproximación rompería el contrato: el matcher descartaría
starts válidos $\Rightarrow$ matches reales perdidos. Eso hace que
el algoritmo sea inviable.

\section{Costo en mmap-spilled}

El \texttt{IndexBacking::heap\_slice\_mut} no funciona sobre
mmap-spilled (es read-only). Para el caso mmap, el motor:

\begin{enumerate}
    \item Abre el archivo en modo \texttt{MmapMut} durante el
    delta update.
    \item Aplica las mutaciones.
    \item \texttt{flush()} para sincronizar a disco.
    \item Vuelve a \texttt{Mmap} read-only.
\end{enumerate}

\noindent
Costo del remap: $\sim$1--3 ms en Linux/Windows. Para una sola
arista delta no compensa; para batches de muchas aristas (decenas
de inserts), la amortización funciona.

\section{Sincronización readers/writer}

El \texttt{KHopReach} expone \texttt{min/max\_first} bajo
\texttt{RwLock} (capítulo~\ref{cap:khop-detalle}). El delta update
toma write-lock mientras muta; los matchers concurrentes ven el
estado consistente vía read-lock. Bloquea, pero la mutación es
\(O(\)visited\()\) que es chico ---típico < 1 ms---.

Para escenarios de altísima concurrencia, una versión wait-free
sería deseable (similar al \texttt{ArcSwap} del streaming backend,
capítulo~\ref{cap:streaming-mirror}); por ahora, \texttt{RwLock}
alcanza.

\section{Tests de equivalencia}

Antes de shipping, dos tests críticos:

\begin{enumerate}
    \item \textbf{Property test}: para un grafo aleatorio + secuencia
    aleatoria de \texttt{add\_relation}, comparar
    \texttt{KHopReach::build} (rebuild from scratch) vs
    \texttt{delta\_update} aplicado uno a uno. Iguales para todos los
    vértices (excepto cuando el cascade budget aborta y el delta es
    sobre-aproximación más laxa que el rebuild).
    \item \textbf{Soundness test}: ningún matcher hot path debe
    podar un start válido después de delta updates. Random hunts
    + verificación contra DFS sin podas.
\end{enumerate}

\section{Cuándo el delta supera al rebuild}

Estimación rough:

\begin{itemize}
    \item Rebuild: \(O(|\Lc_V| \cdot \texttt{MAX\_DEPTH} \cdot
    |\E|)\). Para \(|\E| = 10^7\): $\sim$10 segundos.
    \item Delta: \(O(\texttt{MAX\_DEPTH} \cdot d_{\mathrm{reverse}})\)
    por arista. Para \(d_{\mathrm{reverse}} = 10\) y
    \texttt{MAX\_DEPTH = 4}: $\sim$40 lookups por arista. Para 1000
    aristas insertadas: $\sim$40\,000 ops, $\sim$1 ms.
\end{itemize}

\noindent
Para grafos típicos, el delta es \(\sim 10^4\times\) más rápido que
el rebuild. Para hubs patológicos, el cascade budget detecta y se
cae al rebuild ---peor caso es un poco peor que rebuild puro,
porque pagamos el costo abortado más el rebuild.

\section{Cronograma estimado de implementación}

Stage 2b (delta khop\_reach): $\sim$2 semanas. La propagación
inversa es bien-comprendida; el grueso del trabajo es el property
test de equivalencia delta-vs-rebuild.

\section{Resumen}

\begin{itemize}
    \item Delta WCOJ Propagation: actualiza incremental con
    propagación reverse acotada por \texttt{MAX\_DEPTH}.
    \item Round 1: efecto directo en \(u\). Rounds 2..MAX\_DEPTH:
    cascada vía reverse adjacency, abortar rama si delta vacío.
    \item Cascade budget: si la cascada excede umbral
    (default 100K), invalidar y dejar rebuild.
    \item Soundness por monotonicidad: solo agregamos bits, sólo
    bajamos min y subimos max. Sub-aproximación es imposible.
    \item Costo amortizado: \(O(\texttt{MAX\_DEPTH} \cdot
    d_{\mathrm{reverse}})\) por arista; mucho mejor que rebuild
    \(O(|\E|)\).
    \item Estimado de implementación: $\sim$2 semanas (stage 2b
    futuro).
\end{itemize}

\begin{ejercicios}
\begin{enumerate}
    \item \starred{} Para un grafo donde una nueva arista
    \(u \to v\) tiene \(\bits[v] = \texttt{0b001000}\) (solo bit
    IP) y \(\bits[u] = \texttt{0b000001}\) (solo bit Domain),
    calculá \(\Delta_u\) y los bits resultantes en \(\bits[u]\).

    \item ¿Por qué \texttt{visited} se trackea separado de
    \texttt{frontier}? ¿Qué pasaría sin él?

    \item \starred{} Demostrá que aplicar
    \texttt{delta\_update} \(N\) veces en orden produce
    el mismo \(\bits\) que un rebuild después de las \(N\)
    aristas. (Pista: inducción + monotonicidad.)

    \item Diseñá un grafo adversarial donde el cascade budget
    aborta. ¿Cuál es la cota de visited? ¿Qué pasaría sin el
    budget?

    \item Si el motor implementara \emph{undo} de aristas
    (DELETE), ¿la propagación delta seguiría funcionando? ¿Qué se
    rompe?
\end{enumerate}
\end{ejercicios}
